% !PRINTED 461
\weberpage{461}

Setzen wir für $t$ in $\chi(t)$ die verschiedenen Grössen (4), so
erhalten wir $n$ Grössen, eine aus jedem der conjugirten Körper
\weq{6}{\Theta = \chi(\alpha),\ \Theta_1 = \chi(\alpha_1) \ldots \Theta_{n-1} = \chi(\alpha_{n-1}),}
und diese heissen die mit $\Theta$ \gesperrt{conjugirten Grössen}.

Jede symmetrische Function dieser Grössen ist zugleich eine
symmetrische Function der Wurzeln der Gleichung (3), und mithin in
$\Omega$ enthalten.

Unter diesen symmetrischen Functionen wollen wir zwei, die
häufig vorkommen, durch besondere Namen und Bezeichnungen
hervorheben; es ist die Summe
\weq{7}{S(\Theta) = \Theta + \Theta_1 + \Theta_2 + \cdots + \Theta_{n-1},}
die wir die \gesperrt{Spur} von $\Theta$ nennen, und das Product
\weq{8}{N(\Theta) = \Theta \Theta_1 \Theta_2 \ldots \Theta_{n-1},}
das die \gesperrt{Norm} von $\Theta$ heisst. Conjugirte Zahlen haben hiernach
dieselben Spuren und Normen.

Das Product
\weq{9}{(t - \Theta)(t - \Theta_1) \ldots (t - \Theta_{n-1}) = \Phi(t)}
ist eine ganze rationale Function $n^{\mathrm{ten}}$ Grades von $t$ in
$\Omega$ und ihre Wurzeln sind $\Theta$ und die mit $\Theta$
conjugirten Grössen. Daraus ergiebt sich der Satz:

\gesperrt{1. Jede Grösse in $\Omega(\alpha)$ ist Wurzel einer
Gleichung $n^{\mathrm{ten}}$ Grades in $\Omega$, deren übrige Wurzeln
die mit $\Theta$ conjugirten Grössen sind.}

Die Berechnung der Function $\Phi(t)$ aus $f(x)$ ist nichts
Anderes als die von einem anderen Gesichtspunkte im vierten und
sechsten Abschnitt betrachtete \gesperrt{Tschirnhausen-Transformation}.

Wir haben nun die Function $\Phi(t)$ auf ihre Irreducibilität zu
untersuchen.

Wenn die Function $\Phi(t)$ reducibel ist, so hat sie einen
irreducibeln Factor $\varphi(t)$, in dem wir den Co\"efficienten der
höchsten Potenz von $t$ gleich 1 annehmen können, und jeder solche
Factor verschwindet wenigstens für einen der Werthe $\Theta,
\Theta_1 \ldots \Theta_{n-1}$.

Es sei also
\[
\varphi(\Theta) = \varphi[\chi(\alpha)] = 0,
\]
d.\,h. die Gleichungen $\varphi[\chi(x)] = 0$ und $f(x) = 0$ haben
eine gemeinsame Wurzel. Da aber $f(x)$ irreducibel vorausgesetzt
ist, so muss $\varphi[\chi(x)]$ nach dem Theorem II, \S.\,141 für alle Wurzeln
%CONT>>
